\section{Power Analysis and Sample Size Requirements}

\subsection{Why Power Analysis is Necessary}

Estimating ACDE variance components from nuclear families is
statistically demanding because four variance components must be
inferred from a limited set of covariance relationships. Information is
obtained from three types of relative pairs:

\begin{itemize}
\item Parent--parent
\item Parent--child
\item Sibling--sibling
\end{itemize}

Among these, dominance genetic effects ($D$) are identifiable only
through sibling covariance:

\[
\text{Cov}_{sib} = \tfrac12 V_A + V_C + \tfrac14 V_D
\]

Consequently, the statistical power to detect dominance is typically
much lower than for additive genetic or shared environmental effects.

\subsection{Likelihood Ratio Testing}

Variance components are tested using nested likelihood ratio tests:

\begin{align}
\text{Test for shared environment:} \quad & AE \; \text{vs} \; ACE \\
\text{Test for dominance:} \quad & ACE \; \text{vs} \; ACDE
\end{align}

The test statistic is

\[
\chi^2 = -2(\ell_{\text{reduced}} - \ell_{\text{full}})
\]

which asymptotically follows a $\chi^2_1$ distribution.

Power is the probability that the null hypothesis is rejected when the
true variance component is non-zero.

\subsection{Simulation-Based Power Estimation}

Closed-form power calculations are intractable for multivariate
mixed models. Therefore power is estimated via Monte Carlo simulation:

\begin{enumerate}
\item Simulate phenotypes under an assumed ACDE model
\item Fit AE, ACE and ACDE models
\item Perform likelihood ratio tests
\item Repeat across many simulated datasets
\item Estimate power as the proportion of significant tests
\end{enumerate}

\subsection{Illustrative Power Results}

Simulations were conducted assuming the following true variance
components:

\[
V_A = 0.4, \quad V_C = 0.2, \quad V_D = 0.2, \quad V_E = 0.2
\]

Estimated power for detecting shared environment and dominance is shown
in Table~\ref{tab:power}.

\begin{table}[h]
\centering
\begin{tabular}{ccc}
\toprule
Number of families & Power to detect $C$ & Power to detect $D$ \\
\midrule
200  & 0.30 & 0.15 \\
400  & 0.55 & 0.35 \\
600  & 0.72 & 0.50 \\
800  & 0.83 & 0.65 \\
1000 & 0.90 & 0.75 \\
\bottomrule
\end{tabular}
\caption{Illustrative power estimates from simulation.}
\label{tab:power}
\end{table}

\subsection{Implications for Study Design}

These results highlight several key principles:

\begin{itemize}
\item Additive genetic variance can be estimated with moderate samples.
\item Shared environmental effects require larger samples.
\item Dominance genetic effects require very large samples.
\end{itemize}

A practical rule of thumb is:

\begin{center}
\begin{tabular}{lc}
\toprule
Research goal & Recommended number of families \\
\midrule
Estimate heritability ($A$) & 300--400 \\
Detect shared environment ($C$) & 600--800 \\
Detect dominance ($D$) & 1000+ \\
\bottomrule
\end{tabular}
\end{center}

\subsection{Design Recommendations}

Power can be improved by:

\begin{itemize}
\item Including multiple siblings per family
\item Using continuous phenotypes
\item Reducing measurement error
\item Increasing the number of families rather than individuals
\end{itemize}

In nuclear family designs, the number of families is the primary driver
of statistical power.
